\chapter{System Testing and Evaluation}

The testing and evaluation phase verified that the Intelligent Compliance Management System (ICMS) meets the functional requirements and architectural specifications defined in the design phase. The process ensures that the platform provides security governance, accurate telemetry ingestion, and a stable user experience. Assessment methodologies included Graphical User Interface (GUI) validation, Usability assessments, Performance benchmarking, Exception Handling verification, and Security testing. These tests identified edge cases and validated the deterministic compliance mapping engine.

\section{Critical Appraisal and Limitations}
The ICMS platform demonstrates technical functionality in mapping heterogeneous Security Configuration Assessment (SCA) data to structured GRC controls. The asynchronous processing tier, using Celery and Redis, decouples SIEM polling from the user interface to maintain responsiveness during data ingestion. However, an engineering appraisal reveals limitations in the architecture. Periodic polling of the Wazuh REST API introduces latency between endpoint configuration changes and their reflection on the dashboard. This is suitable for periodic auditing but does not provide the near-real-time synchronization possible with a socket-based streaming architecture. Additionally, generating multi-page PDF audit reports is CPU-intensive; while managed by background workers, it increases the load on the containerized environment during reporting periods.

\section{Testing Environment and Methodology}
The Quality Assurance (QA) process provided full-stack coverage of the ICMS platform. Muhammad Hammad Chaghtai validated the frontend and integration layers, testing the GUI and usability of the React SPA. This included verifying route guarding, state persistence in the \texttt{AuthContext}, theme toggling, and the performance of Recharts-based visualizations. Because the Wazuh SIEM manager was hosted on localized hardware, Muhammad Ali Bukhari validated the backend telemetry modules. This process included testing the Python client's authentication with the SIEM node, verifying network exception handling, and ensuring the integrity of API synchronization tasks. Testing occurred in a distributed containerized environment that simulated standard corporate home-lab conditions.

\section{Functional System Test Cases}
The following fourteen test cases constitute the QA suite used to validate the ICMS platform across IAM, frontend UI, backend API, SIEM telemetry, and reporting layers. Positive and negative scenarios are documented to demonstrate system correctness and exception handling.

% ===================== TC-01 =====================
\begin{table}[H]
    \centering
    \caption{TC-01: Valid Superadmin Authentication (Positive)}
    \label{tab:tc_01}
    \renewcommand{\arraystretch}{1.4}
    \resizebox{\textwidth}{!}{
    \begin{tabular}{|>{\raggedright\arraybackslash\bfseries}p{4.5cm}|>{\raggedright\arraybackslash}p{11.5cm}|}
        \hline
        \rowcolor{blue!10}
        Module Name & Identity and Access Management \\ \hline
        Test Case ID & TC-01 \\ \hline
        Use Case ID & UC-01 \\ \hline
        Tester Name(s) & Muhammad Hammad Chaghtai \\ \hline
        Test Description & Verifying that a Superadmin user can successfully authenticate against the ICMS platform using valid credentials and is navigated to the correct role-specific dashboard upon token issuance. \\ \hline
        Pre-Requisites & 1. A Superadmin account exists in the PostgreSQL \texttt{UserProfile} table. \newline 2. The Django backend API is active on port 8000. \\ \hline
        Environment & React SPA (Windows Client, Chrome 134), Django/PostgreSQL (Docker Ubuntu) \\ \hline
        Test Scenario & A Superadmin enters correct email and password credentials on the login screen. \\ \hline
        Test Steps & 1. Navigate to \url{http://localhost:5173/login}. \newline 2. Enter credentials: \texttt{superadmin@icms.local} and the correct password. \newline 3. Click "Login" and observe routing behavior and localStorage state. \\ \hline
        Test Inputs & Email: \texttt{superadmin@icms.local}, Password: \texttt{[valid credential]} \\ \hline
        Expected Results & The backend should issue a JWT access token and refresh token. The \texttt{AuthContext} should store these tokens, and the React Router should redirect to \texttt{/superadmin-dashboard}. \\ \hline
        Actual Results & HTTP 200 received from \texttt{/api/auth/login/}; JWT stored in \texttt{localStorage} under \texttt{grc\_access\_token}; React Router pushed to \texttt{/superadmin-dashboard} in 120ms. \\ \hline
        Status & \textcolor{green!60!black}{\textbf{PASS}} \\ \hline
        Comments & Token issuance and client-side routing were completed well within acceptable latency thresholds for a LAN-connected environment. \\ \hline
    \end{tabular}
    }
\end{table}

% ===================== TC-02 =====================
\begin{table}[H]
    \centering
    \caption{TC-02: Invalid Credential Rejection (Negative)}
    \label{tab:tc_02}
    \renewcommand{\arraystretch}{1.4}
    \resizebox{\textwidth}{!}{
    \begin{tabular}{|>{\raggedright\arraybackslash\bfseries}p{4.5cm}|>{\raggedright\arraybackslash}p{11.5cm}|}
        \hline
        \rowcolor{blue!10}
        Module Name & Identity and Access Management \\ \hline
        Test Case ID & TC-02 \\ \hline
        Use Case ID & UC-01 \\ \hline
        Tester Name(s) & Muhammad Hammad Chaghtai \\ \hline
        Test Description & Verifying that the authentication endpoint rejects login attempts containing an incorrect password without issuing any session tokens, and that the UI communicates this failure visually to the user. \\ \hline
        Pre-Requisites & 1. A valid user account exists in the database. \newline 2. The login form is rendered and interactive. \\ \hline
        Environment & React SPA (Windows Client, Chrome 134), Django/PostgreSQL (Docker Ubuntu) \\ \hline
        Test Scenario & A user deliberately enters a wrong password to confirm the system does not grant unauthorized access. \\ \hline
        Test Steps & 1. Navigate to \url{http://localhost:5173/login}. \newline 2. Enter \texttt{superadmin@icms.local} with an intentionally incorrect password. \newline 3. Click "Login" and observe the HTTP response and UI feedback. \\ \hline
        Test Inputs & Email: \texttt{superadmin@icms.local}, Password: \texttt{wrongpassword123} \\ \hline
        Expected Results & The backend should return HTTP 401 Unauthorized. No JWT should be stored in \texttt{localStorage}. The UI should display a visible error notification. \\ \hline
        Actual Results & HTTP 401 received from Django; no tokens stored; the \texttt{Toast} notification component rendered a red alert: "Invalid credentials. Please try again." within 85ms. \\ \hline
        Status & \textcolor{green!60!black}{\textbf{PASS}} \\ \hline
        Comments & The rejection path is handled entirely on the backend; no client-side credential caching occurred, confirming the security of the session management logic. \\ \hline
    \end{tabular}
    }
\end{table}

% ===================== TC-03 =====================
\begin{table}[H]
    \centering
    \caption{TC-03: Global Theme Toggling — Light/Dark Mode (Positive)}
    \label{tab:tc_03}
    \renewcommand{\arraystretch}{1.4}
    \resizebox{\textwidth}{!}{
    \begin{tabular}{|>{\raggedright\arraybackslash\bfseries}p{4.5cm}|>{\raggedright\arraybackslash}p{11.5cm}|}
        \hline
        \rowcolor{blue!10}
        Module Name & Frontend State and Visualization Layer \\ \hline
        Test Case ID & TC-03 \\ \hline
        Use Case ID & UC-04 \\ \hline
        Tester Name(s) & Muhammad Hammad Chaghtai \\ \hline
        Test Description & Verifying that the global \texttt{ThemeContext} provider correctly propagates a theme state change across all mounted React components, switching the UI between Light and Dark palettes without requiring a page reload. \\ \hline
        Pre-Requisites & 1. User is authenticated and viewing the main dashboard. \newline 2. The theme toggle control is visible in the navigation bar. \\ \hline
        Environment & React SPA (Windows Client, Chrome 134) \\ \hline
        Test Scenario & A logged-in administrator clicks the theme toggle button while on the dashboard. \\ \hline
        Test Steps & 1. Log in and confirm the default Light Mode is active. \newline 2. Click the sun/moon icon toggle in the top navigation bar. \newline 3. Inspect the DOM's root element for the applied CSS class and verify all component colors update. \\ \hline
        Test Inputs & Initial Theme: \texttt{light}, Toggle Action: \texttt{onClick} \\ \hline
        Expected Results & The \texttt{ThemeContext} state should update from \texttt{light} to \texttt{dark}. All components subscribing to the context should re-render with the Dark Mode CSS variable palette, with no full-page reload. \\ \hline
        Actual Results & \texttt{ThemeContext} state updated immediately; the \texttt{data-theme="dark"} attribute applied to the root \texttt{<div>}; all sidebar, card, and table components re-rendered with the dark palette in under 30ms. \\ \hline
        Status & \textcolor{green!60!black}{\textbf{PASS}} \\ \hline
        Comments & The context-based theme architecture successfully avoids prop drilling and ensures consistent theming across all 20+ React components on the dashboard. \\ \hline
    \end{tabular}
    }
\end{table}

% ===================== TC-04 =====================
\begin{table}[H]
    \centering
    \caption{TC-04: UI RBAC View Segregation (Positive)}
    \label{tab:tc_04}
    \renewcommand{\arraystretch}{1.4}
    \resizebox{\textwidth}{!}{
    \begin{tabular}{|>{\raggedright\arraybackslash\bfseries}p{4.5cm}|>{\raggedright\arraybackslash}p{11.5cm}|}
        \hline
        \rowcolor{blue!10}
        Module Name & Identity and Access Management \\ \hline
        Test Case ID & TC-04 \\ \hline
        Use Case ID & UC-05 \\ \hline
        Tester Name(s) & Muhammad Hammad Chaghtai \\ \hline
        Test Description & Verifying that the React frontend enforces role-based component rendering by conditionally omitting privileged navigation items from the DOM when an Auditor-level session is active. \\ \hline
        Pre-Requisites & 1. Two accounts are available: one Auditor, one Superadmin. \newline 2. Browser DevTools are open to inspect the DOM. \\ \hline
        Environment & React SPA (Windows Client, Chrome 134), Django/PostgreSQL (Docker Ubuntu) \\ \hline
        Test Scenario & An Auditor logs in and inspects the sidebar navigation to confirm restricted items are absent from the DOM. \\ \hline
        Test Steps & 1. Authenticate as an Auditor. \newline 2. Observe the sidebar navigation menu after the dashboard loads. \newline 3. Use Chrome DevTools to verify the "Manage Admins" and "Mapping Configuration" elements are not rendered in the DOM. \\ \hline
        Test Inputs & Role: \texttt{auditor}, JWT Role Claim: \texttt{\{"role": "auditor"\}} \\ \hline
        Expected Results & The \texttt{Sidebar} component should conditionally omit the "Manage Admins" and "Mapping Configuration" navigation items. These routes should also redirect to a 403 page if accessed directly. \\ \hline
        Actual Results & Sidebar rendered 4 navigation items for Auditor vs. 9 for Superadmin. DevTools confirmed the privileged \texttt{<NavItem>} components were not mounted in the DOM, not merely hidden with CSS. \\ \hline
        Status & \textcolor{green!60!black}{\textbf{PASS}} \\ \hline
        Comments & DOM-level omission (as opposed to CSS \texttt{display:none}) is the correct security pattern as it prevents advanced users from unhiding elements via browser console commands. \\ \hline
    \end{tabular}
    }
\end{table}

% ===================== TC-05 =====================
\begin{table}[H]
    \centering
    \caption{TC-05: Unauthorized API Write Attempt (Negative)}
    \label{tab:tc_05}
    \renewcommand{\arraystretch}{1.4}
    \resizebox{\textwidth}{!}{
    \begin{tabular}{|>{\raggedright\arraybackslash\bfseries}p{4.5cm}|>{\raggedright\arraybackslash}p{11.5cm}|}
        \hline
        \rowcolor{blue!10}
        Module Name & Identity and Access Management \\ \hline
        Test Case ID & TC-05 \\ \hline
        Use Case ID & UC-05 \\ \hline
        Tester Name(s) & Muhammad Hammad Chaghtai \& Muhammad Ali Bukhari \\ \hline
        Test Description & Verifying that the Django REST Framework permission layer blocks write operations originating from an Auditor-level JWT, even when the HTTP request is constructed manually to bypass the UI restrictions. \\ \hline
        Pre-Requisites & 1. User is authenticated as an Auditor. \newline 2. Postman or curl is used to craft a direct API request. \\ \hline
        Environment & Postman (Windows Client), Django/PostgreSQL (Docker Ubuntu) \\ \hline
        Test Scenario & An Auditor uses Postman to construct a direct POST request to the framework management endpoint, attempting to create a new compliance control mapping. \\ \hline
        Test Steps & 1. Authenticate as Auditor via \texttt{/api/auth/login/} to receive a restricted JWT. \newline 2. Attach the JWT and construct a \texttt{POST} request to \texttt{/api/frameworks/}. \newline 3. Submit the request and record the HTTP response code and body. \\ \hline
        Test Inputs & Role: \texttt{auditor}, Endpoint: \texttt{POST /api/frameworks/}, Payload: \texttt{\{"name": "Unauthorized\_Frame"\}} \\ \hline
        Expected Results & DRF \texttt{IsAdminOrSuperAdmin} permission class should intercept the request and return HTTP 403 Forbidden. No database record should be created. \\ \hline
        Actual Results & HTTP 403 Forbidden returned in 52ms; response body: \texttt{\{"detail": "You do not have permission to perform this action."\}}; zero database writes confirmed via query log. \\ \hline
        Status & \textcolor{green!60!black}{\textbf{PASS}} \\ \hline
        Comments & The permission layer operates at the view level, independent of the UI, ensuring API security even against direct HTTP attacks. \\ \hline
    \end{tabular}
    }
\end{table}

% ===================== TC-06 =====================
\begin{table}[H]
    \centering
    \caption{TC-06: Deterministic Framework Mapping (Positive)}
    \label{tab:tc_06}
    \renewcommand{\arraystretch}{1.4}
    \resizebox{\textwidth}{!}{
    \begin{tabular}{|>{\raggedright\arraybackslash\bfseries}p{4.5cm}|>{\raggedright\arraybackslash}p{11.5cm}|}
        \hline
        \rowcolor{blue!10}
        Module Name & Governance and Compliance Data Engine \\ \hline
        Test Case ID & TC-06 \\ \hline
        Use Case ID & UC-03 \\ \hline
        Tester Name(s) & Muhammad Hammad Chaghtai \\ \hline
        Test Description & Verifying that a Superadmin can create a valid regulatory mapping linking a specific Wazuh SCA Rule ID to a named ISO 27001 control, and that the record is persisted correctly in the database. \\ \hline
        Pre-Requisites & 1. ISO 27001 Framework and Control \texttt{A.12.1.2} are seeded in the database. \newline 2. Wazuh Rule ID \texttt{19502} does not yet have a mapping. \\ \hline
        Environment & React SPA (Windows Client, Chrome 134), Django/PostgreSQL (Docker Ubuntu) \\ \hline
        Test Scenario & A Superadmin navigates to the mapping panel and creates a new rule-to-control entry. \\ \hline
        Test Steps & 1. Log in as Superadmin and navigate to "Mapping Configuration". \newline 2. Input Wazuh Rule ID \texttt{19502}, select target control \texttt{A.12.1.2 (Change Management)}, and click "Save". \newline 3. Verify the new row appears in the mapping table on the page. \\ \hline
        Test Inputs & Wazuh Rule ID: \texttt{19502}, Control: \texttt{A.12.1.2 — Change Management}, Framework: \texttt{ISO 27001} \\ \hline
        Expected Results & A new \texttt{WazuhMapping} record should be persisted to the PostgreSQL database and reflected in the UI mapping table. \\ \hline
        Actual Results & \texttt{POST /api/mappings/} returned HTTP 201 Created; database commit confirmed in 110ms; the new mapping appeared in the UI table immediately after the Axios response resolved. \\ \hline
        Status & \textcolor{green!60!black}{\textbf{PASS}} \\ \hline
        Comments & The successful creation of this mapping is foundational to the entire automated scoring pipeline; the 110ms DB commit time is well within acceptable thresholds. \\ \hline
    \end{tabular}
    }
\end{table}

% ===================== TC-07 =====================
\begin{table}[H]
    \centering
    \caption{TC-07: Duplicate Mapping Collision (Negative)}
    \label{tab:tc_07}
    \renewcommand{\arraystretch}{1.4}
    \resizebox{\textwidth}{!}{
    \begin{tabular}{|>{\raggedright\arraybackslash\bfseries}p{4.5cm}|>{\raggedright\arraybackslash}p{11.5cm}|}
        \hline
        \rowcolor{blue!10}
        Module Name & Governance and Compliance Data Engine \\ \hline
        Test Case ID & TC-07 \\ \hline
        Use Case ID & UC-03 \\ \hline
        Tester Name(s) & Muhammad Hammad Chaghtai \\ \hline
        Test Description & Verifying that the database unique constraint on the \texttt{WazuhMapping} model prevents the creation of a duplicate entry for an already-mapped Wazuh Rule ID, and that the backend returns a structured validation error to the UI. \\ \hline
        Pre-Requisites & 1. Wazuh Rule ID \texttt{19502} is already mapped to Control \texttt{A.12.1.2} (established in TC-06). \newline 2. User is authenticated as Superadmin. \\ \hline
        Environment & React SPA (Windows Client, Chrome 134), Django/PostgreSQL (Docker Ubuntu) \\ \hline
        Test Scenario & A Superadmin attempts to create an identical mapping entry for the same Rule ID and Control combination. \\ \hline
        Test Steps & 1. Navigate to "Mapping Configuration". \newline 2. Attempt to save a new mapping using the same Rule ID \texttt{19502} and the same Control \texttt{A.12.1.2}. \newline 3. Observe the error response from the backend. \\ \hline
        Test Inputs & Wazuh Rule ID: \texttt{19502}, Control: \texttt{A.12.1.2}, Framework: \texttt{ISO 27001} (duplicate) \\ \hline
        Expected Results & PostgreSQL should raise an \texttt{IntegrityError} for the unique-together constraint. The DRF serializer should catch this and return HTTP 400 Bad Request. The UI should display a distinct error alert. \\ \hline
        Actual Results & \texttt{IntegrityError} caught by the DRF serializer; HTTP 400 returned with body \texttt{\{"non\_field\_errors": ["Duplicate Mapping Exception: This rule is already mapped."]\}}; UI displayed an orange warning toast in 95ms. \\ \hline
        Status & \textcolor{green!60!black}{\textbf{PASS}} \\ \hline
        Comments & The constraint is enforced at the database level, making it impossible to create duplicate mappings even via direct API calls. \\ \hline
    \end{tabular}
    }
\end{table}

% ===================== TC-08 =====================
\begin{table}[H]
    \centering
    \caption{TC-08: Wazuh API Handshake and Token Retrieval (Positive)}
    \label{tab:tc_08}
    \renewcommand{\arraystretch}{1.4}
    \resizebox{\textwidth}{!}{
    \begin{tabular}{|>{\raggedright\arraybackslash\bfseries}p{4.5cm}|>{\raggedright\arraybackslash}p{11.5cm}|}
        \hline
        \rowcolor{blue!10}
        Module Name & SIEM Telemetry and Synchronization \\ \hline
        Test Case ID & TC-08 \\ \hline
        Use Case ID & UC-02 \\ \hline
        Tester Name(s) & Muhammad Ali Bukhari \\ \hline
        Test Description & Validating the \texttt{WazuhAPIClient} service's capability to establish an authenticated session with the Wazuh Manager REST API by exchanging Basic Auth credentials for a short-lived Bearer token. \\ \hline
        Pre-Requisites & 1. Wazuh Manager is active at \texttt{192.168.100.100:55000}. \newline 2. Environment variables \texttt{WAZUH\_API\_USER} and \texttt{WAZUH\_API\_PASSWORD} are configured. \\ \hline
        Environment & Django Backend (Docker Ubuntu), Wazuh Manager 4.x (Physical Laptop) \\ \hline
        Test Scenario & The ICMS backend initializes the \texttt{WazuhAPIClient} and calls the \texttt{authenticate()} method. \\ \hline
        Test Steps & 1. Trigger \texttt{WazuhAPIClient.authenticate()} from the Django shell. \newline 2. Monitor outbound HTTPS traffic on Port 55000. \newline 3. Confirm the Bearer token is stored on the client instance. \\ \hline
        Test Inputs & Host: \texttt{192.168.100.100}, Port: \texttt{55000}, SSL Verify: \texttt{False} \\ \hline
        Expected Results & The client should call \texttt{GET /security/user/authenticate} with Basic Auth and store the resulting JWT as \texttt{self.\_token}. \\ \hline
        Actual Results & \texttt{InsecureRequestWarning} suppressed; HTTP 200 received; Bearer token retrieved and stored in 0.8s despite self-signed certificate overhead. \\ \hline
        Status & \textcolor{green!60!black}{\textbf{PASS}} \\ \hline
        Comments & Disabling SSL verification is an accepted trade-off in a home-lab deployment; production environments should use a valid CA-signed certificate. \\ \hline
    \end{tabular}
    }
\end{table}

% ===================== TC-09 =====================
\begin{table}[H]
    \centering
    \caption{TC-09: Wazuh Connection Timeout Exception (Negative)}
    \label{tab:tc_09}
    \renewcommand{\arraystretch}{1.4}
    \resizebox{\textwidth}{!}{
    \begin{tabular}{|>{\raggedright\arraybackslash\bfseries}p{4.5cm}|>{\raggedright\arraybackslash}p{11.5cm}|}
        \hline
        \rowcolor{blue!10}
        Module Name & SIEM Telemetry and Synchronization \\ \hline
        Test Case ID & TC-09 \\ \hline
        Use Case ID & UC-02 \\ \hline
        Tester Name(s) & Muhammad Ali Bukhari \\ \hline
        Test Description & Simulating a complete network failure during SIEM synchronization to validate that the Celery worker catches the timeout exception and triggers an exponential backoff retry rather than crashing silently. \\ \hline
        Pre-Requisites & 1. Wazuh Manager taken offline (firewall rule blocks Port 55000). \newline 2. Celery worker configured with \texttt{max\_retries=3}. \\ \hline
        Environment & Django Backend (Docker Ubuntu), Wazuh Manager 4.x (Offline) \\ \hline
        Test Scenario & A scheduled telemetry synchronization task fires while the SIEM node is unreachable. \\ \hline
        Test Steps & 1. Block outbound traffic to Port 55000. \newline 2. Execute \texttt{sync\_wazuh\_scans.delay()}. \newline 3. Monitor Celery logs for exception handling and retry behavior. \\ \hline
        Test Inputs & Target Host: \texttt{192.168.100.100:55000}, Network State: \texttt{Connection Refused} \\ \hline
        Expected Results & \texttt{requests.exceptions.Timeout} caught; Celery calls \texttt{self.retry(countdown=60)} to re-enqueue the task. \\ \hline
        Actual Results & Timeout raised after 15s; Celery logged "Task retrying in 60s"; task completed successfully on second attempt after connectivity was restored. \\ \hline
        Status & \textcolor{orange}{\textbf{PASS (Handled Exception)}} \\ \hline
        Comments & The retry mechanism prevents data loss during transient network failures; the 60-second countdown avoids flooding a recovering SIEM node. \\ \hline
    \end{tabular}
    }
\end{table}

% ===================== TC-10 =====================
\begin{table}[H]
    \centering
    \caption{TC-10: Asynchronous Telemetry Polling via Celery Beat (Positive)}
    \label{tab:tc_10}
    \renewcommand{\arraystretch}{1.4}
    \resizebox{\textwidth}{!}{
    \begin{tabular}{|>{\raggedright\arraybackslash\bfseries}p{4.5cm}|>{\raggedright\arraybackslash}p{11.5cm}|}
        \hline
        \rowcolor{blue!10}
        Module Name & Asynchronous Processing Engine \\ \hline
        Test Case ID & TC-10 \\ \hline
        Use Case ID & UC-02 \\ \hline
        Tester Name(s) & Muhammad Ali Bukhari \\ \hline
        Test Description & Validating the full end-to-end execution of the \texttt{sync\_wazuh\_scans} Celery Beat periodic task, from queue dispatch to database commit, confirming the API thread is never blocked during the operation. \\ \hline
        Pre-Requisites & 1. Redis broker is active. \newline 2. At least 5 agents registered on Wazuh Manager. \newline 3. A periodic task entry exists in the \texttt{django\_celery\_beat} schedule table. \\ \hline
        Environment & Celery Worker / Redis (Docker Ubuntu), Wazuh Manager 4.x \\ \hline
        Test Scenario & The scheduled synchronization task is triggered by Celery Beat to poll all 5 registered endpoints. \\ \hline
        Test Steps & 1. Monitor Celery Beat logs for the dispatch event. \newline 2. Confirm all 5 agents are polled sequentially. \newline 3. Query the \texttt{ScanResult} table for new records after completion. \\ \hline
        Test Inputs & Task: \texttt{sync\_wazuh\_scans}, Agent Count: \texttt{5}, Trigger: \texttt{Manual} \\ \hline
        Expected Results & All 5 endpoints polled; new \texttt{ScanResult} records committed; Django API process remains responsive throughout. \\ \hline
        Actual Results & 1,200 SCA check rows updated across 5 agents in 12.4s; simultaneous \texttt{GET /api/dashboard/} request responded in 78ms, confirming non-blocking operation. \\ \hline
        Status & \textcolor{green!60!black}{\textbf{PASS}} \\ \hline
        Comments & The 78ms concurrent API response confirms the Celery worker tier decoupling is functioning correctly under active load. \\ \hline
    \end{tabular}
    }
\end{table}

% ===================== TC-11 =====================
\begin{table}[H]
    \centering
    \caption{TC-11: Unstructured JSON Log Ingestion (Positive)}
    \label{tab:tc_11}
    \renewcommand{\arraystretch}{1.4}
    \resizebox{\textwidth}{!}{
    \begin{tabular}{|>{\raggedright\arraybackslash\bfseries}p{4.5cm}|>{\raggedright\arraybackslash}p{11.5cm}|}
        \hline
        \rowcolor{blue!10}
        Module Name & Governance and Compliance Data Engine \\ \hline
        Test Case ID & TC-11 \\ \hline
        Use Case ID & UC-02 \\ \hline
        Tester Name(s) & Muhammad Ali Bukhari \\ \hline
        Test Description & Verifying the integrity of the raw log ingestion pipeline when saving a large, semi-structured Wazuh SCA JSON payload into the PostgreSQL \texttt{jsonb} field on the \texttt{ScanResult} model. \\ \hline
        Pre-Requisites & 1. Wazuh API returning a high-density SCA check payload. \newline 2. PostgreSQL 18 database active and accepting connections. \\ \hline
        Environment & Django Backend (Docker Ubuntu), PostgreSQL 18 \\ \hline
        Test Scenario & A synchronization task receives an 18KB JSON object from the Wazuh API and saves it to \texttt{raw\_log\_data}. \\ \hline
        Test Steps & 1. Execute the sync task for a high-density agent. \newline 2. Query the \texttt{ScanResult} table for the \texttt{raw\_log\_data} field. \newline 3. Compare character count and JSON key structure with the source SIEM payload. \\ \hline
        Test Inputs & Payload Source: Wazuh SCA endpoint, Payload Size: \texttt{18KB}, Field: \texttt{jsonb} \\ \hline
        Expected Results & The entire JSON object stored intact with no truncation, and fully traversable via Django ORM for audit retrieval. \\ \hline
        Actual Results & 18KB payload committed to \texttt{raw\_log\_data} in 1.1s; character-count comparison confirmed zero data loss; JSON tree fully traversable via ORM. \\ \hline
        Status & \textcolor{green!60!black}{\textbf{PASS}} \\ \hline
        Comments & The dual-storage approach (structured boolean + raw JSON) preserves original evidence while enabling fast dashboard queries. \\ \hline
    \end{tabular}
    }
\end{table}

% ===================== TC-12 =====================
\begin{table}[H]
    \centering
    \caption{TC-12: Dynamic Compliance Score Recalculation (Positive)}
    \label{tab:tc_12}
    \renewcommand{\arraystretch}{1.4}
    \resizebox{\textwidth}{!}{
    \begin{tabular}{|>{\raggedright\arraybackslash\bfseries}p{4.5cm}|>{\raggedright\arraybackslash}p{11.5cm}|}
        \hline
        \rowcolor{blue!10}
        Module Name & Asynchronous Processing Engine \\ \hline
        Test Case ID & TC-12 \\ \hline
        Use Case ID & UC-04 \\ \hline
        Tester Name(s) & Muhammad Ali Bukhari \\ \hline
        Test Description & Validating the scoring aggregation logic by confirming that a new scan result resolving a previously non-compliant control triggers a recalculation of the global compliance percentage, reflected immediately in the frontend gauge. \\ \hline
        Pre-Requisites & 1. Agent \texttt{004} has a baseline score of 75.0\%. \newline 2. A failed SCA check for a mapped control exists in the database. \\ \hline
        Environment & Celery Worker (Docker Ubuntu), React SPA (Windows Client) \\ \hline
        Test Scenario & A manual synchronization resolves a previously failing control, incrementing the global score. \\ \hline
        Test Steps & 1. Trigger \texttt{sync\_wazuh\_scans.delay(agent\_id=4)} manually. \newline 2. Query \texttt{DashboardSummaryView} for the updated compliance score. \newline 3. Refresh the dashboard and observe the \texttt{ComplianceGauge} animation. \\ \hline
        Test Inputs & Agent ID: \texttt{004}, Control Weight: \texttt{+6.2\%}, Initial Score: \texttt{75.0\%} \\ \hline
        Expected Results & Global score increments from 75.0\% to 81.2\%; the React dashboard re-renders the \texttt{ComplianceGauge} SVG with the new value. \\ \hline
        Actual Results & Score updated to 81.2\% in the database within 2.4s of task completion; gauge SVG re-animated to the new value in 40ms without a page reload. \\ \hline
        Status & \textcolor{green!60!black}{\textbf{PASS}} \\ \hline
        Comments & The real-time gauge update is achieved through the Axios polling mechanism in the \texttt{Dashboard} component, re-fetching data every 30 seconds. \\ \hline
    \end{tabular}
    }
\end{table}

% ===================== TC-13 =====================
\begin{table}[H]
    \centering
    \caption{TC-13: Comprehensive PDF Audit Report Generation (Positive)}
    \label{tab:tc_13}
    \renewcommand{\arraystretch}{1.4}
    \resizebox{\textwidth}{!}{
    \begin{tabular}{|>{\raggedright\arraybackslash\bfseries}p{4.5cm}|>{\raggedright\arraybackslash}p{11.5cm}|}
        \hline
        \rowcolor{blue!10}
        Module Name & Automated Audit and Reporting Module \\ \hline
        Test Case ID & TC-13 \\ \hline
        Use Case ID & UC-06 \\ \hline
        Tester Name(s) & Muhammad Hammad Chaghtai \& Muhammad Ali Bukhari \\ \hline
        Test Description & Testing the server-side PDF rendering pipeline's ability to aggregate multi-agent compliance data, apply a Django HTML template, and produce a structurally sound binary PDF with correct page breaks and table alignment. \\ \hline
        Pre-Requisites & 1. Database contains at least 50 historical \texttt{ScanResult} records across 5 agents. \newline 2. \texttt{xhtml2pdf} is installed and the report HTML template is configured. \\ \hline
        Environment & Django Backend (Docker Ubuntu), React SPA (Windows Client) \\ \hline
        Test Scenario & An Auditor clicks "Generate Full Report" for the ISO 27001 framework on the Reports page. \\ \hline
        Test Steps & 1. Authenticate as Auditor and navigate to the Reports page. \newline 2. Select framework "ISO 27001" and click "Generate Report". \newline 3. Download the returned PDF binary and inspect page breaks and table layout. \\ \hline
        Test Inputs & Framework: \texttt{ISO 27001}, Date Range: \texttt{Last 30 Days}, Agent Count: \texttt{5} \\ \hline
        Expected Results & A multi-page PDF containing a framework summary, per-agent compliance tables, and a gap analysis section, with clean page breaks and no overlapping elements. \\ \hline
        Actual Results & 15-page PDF document generated in 4.2s; all inter-section page breaks executed correctly without table overlap; all data values matched the live database query. \\ \hline
        Status & \textcolor{green!60!black}{\textbf{PASS}} \\ \hline
        Comments & The 4.2s generation time for a 15-page dossier is acceptable; offloading to the worker tier ensured the auditor's browser session remained interactive throughout. \\ \hline
    \end{tabular}
    }
\end{table}

% ===================== TC-14 =====================
\begin{table}[H]
    \centering
    \caption{TC-14: Asynchronous SMTP Alert Dispatch (Positive)}
    \label{tab:tc_14}
    \renewcommand{\arraystretch}{1.4}
    \resizebox{\textwidth}{!}{
    \begin{tabular}{|>{\raggedright\arraybackslash\bfseries}p{4.5cm}|>{\raggedright\arraybackslash}p{11.5cm}|}
        \hline
        \rowcolor{blue!10}
        Module Name & Asynchronous Processing Engine \\ \hline
        Test Case ID & TC-14 \\ \hline
        Use Case ID & UC-05 \\ \hline
        Tester Name(s) & Muhammad Hammad Chaghtai \& Muhammad Ali Bukhari \\ \hline
        Test Description & Validating the automated alert system's capability to detect a compliance score threshold breach and dispatch a structured HTML email notification via the SMTP relay, without blocking the primary ingestion pipeline. \\ \hline
        Pre-Requisites & 1. System settings configured with an alert threshold of 80\%. \newline 2. SMTP credentials pointing to a MailTrap test inbox are active. \\ \hline
        Environment & Celery Worker (Docker Ubuntu), SMTP Relay (MailTrap) \\ \hline
        Test Scenario & A telemetry synchronization ingests data dropping Agent \texttt{004}'s compliance score below the 80\% critical threshold, triggering the alert pipeline. \\ \hline
        Test Steps & 1. Insert a simulated scan result lowering Agent 004's score to 65.0\%. \newline 2. Monitor Celery logs for the \texttt{send\_compliance\_alert} task dispatch. \newline 3. Check the MailTrap inbox for the received email. \\ \hline
        Test Inputs & Agent ID: \texttt{004}, New Score: \texttt{65.0\%}, Configured Threshold: \texttt{80.0\%} \\ \hline
        Expected Results & The scoring module detects the breach, enqueues a \texttt{send\_compliance\_alert} task, and the worker dispatches an HTML email containing the agent name, current score, and framework details. \\ \hline
        Actual Results & Threshold breach detected post-aggregation; \texttt{send\_compliance\_alert} task dispatched in 1.1s; HTML email received in MailTrap containing Agent 004 metadata, score 65.0\%, and a direct link to the agent detail page. \\ \hline
        Status & \textcolor{green!60!black}{\textbf{PASS}} \\ \hline
        Comments & The dedicated SMTP worker ensures email dispatch latency does not delay subsequent compliance ingestion tasks queued behind it. \\ \hline
    \end{tabular}
    }
\end{table}

\section{Non-Functional System Testing}
The evaluation of the ICMS platform extended beyond functional correctness to assess the system's operational constraints, scalability, and overall user experience. This phase of non-functional testing focused on ensuring that the architectural decisions---such as the separation of concerns and the use of an asynchronous task queue---translate into a performant, secure, and usable application for security analysts.

\subsection{Performance and Load Testing}
The performance evaluation confirmed that the ICMS platform maintains high responsiveness under typical operational loads. Benchmarking revealed that standard API requests, such as retrieving framework lists or mapping configurations, consistently resolve in under 150ms. To handle more intensive I/O operations, such as polling large datasets from the Wazuh SIEM or rendering multi-page PDF reports, the system utilized a Celery worker tier. This architectural choice successfully offloads long-running processes from the main request-response cycle, preventing the Nginx and Django gateway from timing out during heavy stress. Even when processing 50 simultaneous synchronization tasks, the system remained stable, completing all operations without degrading the user interface performance.

\begin{table}[H]
    \centering
    \caption{Performance Test Results}
    \label{tab:perf_results}
    \renewcommand{\arraystretch}{1.4}
    \resizebox{\textwidth}{!}{
    \begin{tabular}{|>{\raggedright\arraybackslash}p{4.5cm}|>{\raggedright\arraybackslash}p{4cm}|>{\raggedright\arraybackslash}p{3.5cm}|>{\raggedright\arraybackslash}p{1.5cm}|}
        \hline
        \rowcolor{blue!10}
        \textbf{Test Scenario} & \textbf{Metric Measured} & \textbf{Result} & \textbf{Status} \\ \hline
        React SPA cold start to Login & Screen load time & Under 1.5 seconds & Pass \\ \hline
        Django API framework list fetch & Response latency & Under 150 ms & Pass \\ \hline
        Wazuh SIEM Sync (50 agents) & Celery execution time & Under 15 seconds & Pass \\ \hline
        Compliance gauge re-render & UI update latency & Under 50 ms & Pass \\ \hline
        ISO 27001 PDF generation (15 pages) & Server rendering time & Under 4.5 seconds & Pass \\ \hline
        SMTP alert email dispatch & Queue delivery time & Under 1.5 seconds & Pass \\ \hline
    \end{tabular}
    }
\end{table}

\subsection{Security and Penetration Testing}
A defense-in-depth strategy was implemented to secure the platform against unauthorized access and data manipulation. The system relies on the Django ORM to provide inherent protection against SQL injection attacks by utilizing parameterized queries for all database interactions. Identity management is handled through a stateless JWT architecture, utilizing short-lived access tokens and rotating refresh tokens to minimize the window for potential session hijacking. Additionally, the integration of TOTP-based 2FA provides a cryptographically secure second layer of authentication, ensuring that administrative access remains protected even in the event of primary credential compromise.

\begin{table}[H]
    \centering
    \caption{Security Test Results}
    \label{tab:sec_results}
    \renewcommand{\arraystretch}{1.4}
    \resizebox{\textwidth}{!}{
    \begin{tabular}{|>{\raggedright\arraybackslash}p{4cm}|>{\raggedright\arraybackslash}p{8cm}|>{\raggedright\arraybackslash}p{1.5cm}|}
        \hline
        \rowcolor{blue!10}
        \textbf{Security Control} & \textbf{Test Description} & \textbf{Status} \\ \hline
        JWT Session Enforcement & Attempted access to protected Dashboard without token; React Router successfully redirected to login. & Pass \\ \hline
        Role-Based Access Control & Attempted to POST mapping rules using an Auditor token; DRF returned HTTP 403 Forbidden. & Pass \\ \hline
        SQL Injection Prevention & Submitted malicious SQL payloads in search fields; Django ORM parameterized queries safely. & Pass \\ \hline
        TOTP 2FA Verification & Attempted to bypass the 2FA prompt after password entry; backend withheld the final access token. & Pass \\ \hline
        XSS Sanitization & Injected HTML scripts into framework descriptions; React DOM safely escaped the output. & Pass \\ \hline
    \end{tabular}
    }
\end{table}

\subsection{Usability and Accessibility Testing}
The user interface is designed to manage cognitive load for analysts in high-pressure Security Operations Center (SOC) environments. The implementation includes a global \texttt{ThemeContext} that allows toggling between Light and Dark modes. This gives analysts control over the UI palette based on local lighting conditions to reduce eye strain. The platform organizes compliance data into paginated and filterable tables to ensure that datasets remain manageable without overwhelming the interface.

\begin{table}[H]
    \centering
    \caption{Usability Test Task Completion Results}
    \label{tab:usability_results}
    \renewcommand{\arraystretch}{1.4}
    \resizebox{\textwidth}{!}{
    \begin{tabular}{|>{\raggedright\arraybackslash}p{6cm}|>{\raggedright\arraybackslash}p{3cm}|>{\raggedright\arraybackslash}p{2.5cm}|>{\raggedright\arraybackslash}p{1.5cm}|}
        \hline
        \rowcolor{blue!10}
        \textbf{Task} & \textbf{Avg. Completion Time} & \textbf{Success Rate} & \textbf{Status} \\ \hline
        Log in and complete 2FA verification & 0.8 minutes & 5 / 5 (100\%) & Pass \\ \hline
        Toggle Light/Dark mode context & 0.2 minutes & 5 / 5 (100\%) & Pass \\ \hline
        Map a Wazuh Rule to an ISO Control & 1.5 minutes & 5 / 5 (100\%) & Pass \\ \hline
        Trigger a manual endpoint scan & 1.0 minutes & 5 / 5 (100\%) & Pass \\ \hline
        Generate and download PDF report & 1.2 minutes & 4 / 5 (80\%) & Pass \\ \hline
    \end{tabular}
    }
\end{table}

\subsection{Compatibility and Deployment Testing}
The ICMS platform underwent compatibility testing to ensure a consistent experience across modern web browsers, including Chrome, Firefox, and Edge. The React SPA maintained its layout integrity and functionality across these environments, confirming the reliability of the frontend codebase. Deployment testing validated the containerized architecture by performing multiple teardowns and spin-ups using Docker Compose. By encapsulating the entire stack---including Django, React, Redis, Celery, and PostgreSQL---the deployment remains deterministic across different host environments. This containerized approach ensures that the platform can be reliably deployed on any Ubuntu-based machine without encountering dependency conflicts.

\begin{table}[H]
    \centering
    \caption{Compatibility and Deployment Test Results}
    \label{tab:compat_results}
    \renewcommand{\arraystretch}{1.4}
    \resizebox{\textwidth}{!}{
    \begin{tabular}{|>{\raggedright\arraybackslash}p{4cm}|>{\raggedright\arraybackslash}p{8cm}|>{\raggedright\arraybackslash}p{1.5cm}|}
        \hline
        \rowcolor{blue!10}
        \textbf{Environment / Scenario} & \textbf{Test Description} & \textbf{Status} \\ \hline
        Cross-Browser UI Rendering & Executed React SPA on Chrome, Firefox, and Edge; Flexbox layouts and charts rendered correctly. & Pass \\ \hline
        Docker Compose Teardown & Executed \texttt{docker-compose down -v} and rebuilt stack; all 7 containers initialized with Exit Code 0. & Pass \\ \hline
        Database Volume Persistence & Restarted PostgreSQL container; verified that previous user accounts and compliance data remained intact. & Pass \\ \hline
        Celery Broker Recovery & Temporarily stopped Redis container; verified Celery gracefully resumed task queuing upon Redis restart. & Pass \\ \hline
    \end{tabular}
    }
\end{table}
